\documentclass[conference]{IEEEtran}
\IEEEoverridecommandlockouts

\usepackage{cite}
\usepackage{amsmath,amssymb,amsfonts}
\usepackage{graphicx}
\usepackage{textcomp}
\usepackage{xcolor}
\usepackage{booktabs}
\usepackage{array}

\begin{document}

\title{Constant Bandwidth Server: Critical Evaluation and Transfer\\
{\footnotesize Milestone 3 --- Real-Time Systems Seminar, HSHL}}

\author{\IEEEauthorblockN{Soaib Ferdous}
\IEEEauthorblockA{\textit{BEng in Electronic Engineering} \\
\textit{Hochschule Hamm-Lippstadt}\\
Student ID: 1232368 \\
Group B, Team 6}}

\maketitle

% ============================================================
\section{Concrete Strengths}
% ============================================================

% Annotation: not generic advantages — specific to CBS design choices

\begin{itemize}
  \item \textbf{Formal schedulability guarantee:}
        CBS integrates into EDF analysis as a virtual periodic task with
        $U_s = Q_s/T_s$ $\Rightarrow$ schedulability test is exact,
        not heuristic
  \item \textbf{Temporal isolation:}
        overrunning aperiodic task cannot affect periodic tasks beyond
        $U_s$ bandwidth $\Rightarrow$ stronger than DS, TBS
  \item \textbf{Work-conserving:}
        server uses available slack immediately when job arrives and
        budget is fresh $\Rightarrow$ better aperiodic response time
        than PS
  \item \textbf{Simple schedulability test:}
        same EDF condition as periodic tasks --- no separate analysis
        needed for the server
        \begin{equation}
          U_s + \sum_{i=1}^{n} \frac{C_i}{T_i} \leq 1
        \end{equation}
  \item \textbf{Low runtime overhead:}
        $O(1)$ server state update per event --- feasible on
        resource-constrained embedded hardware
  \item \textbf{Graceful overload handling:}
        overrunning task continues execution at lower EDF priority
        (deadline postponed) $\Rightarrow$ no job is killed or dropped
\end{itemize}

% ============================================================
\section{Concrete Limitations and Risks}
% ============================================================

\begin{itemize}
  \item \textbf{EDF dependency:}
        CBS cannot run on fixed-priority schedulers (RM, DM, AUTOSAR)
        without replacing the scheduler $\Rightarrow$ major deployment
        barrier in industry
  \item \textbf{No hard per-job response time guarantee:}
        CBS guarantees bandwidth, not individual aperiodic job deadlines
        $\Rightarrow$ unsuitable for hard real-time aperiodic tasks
  \item \textbf{Parameter selection is offline and static:}
        $Q_s$ and $T_s$ must be set before runtime. Under highly
        variable aperiodic load, a fixed $(Q_s, T_s)$ pair may be:
        \begin{itemize}
          \item too conservative $\Rightarrow$ wasted bandwidth
          \item too aggressive $\Rightarrow$ insufficient isolation
        \end{itemize}
  \item \textbf{Single-processor formulation:}
        original CBS has no multi-core extension $\Rightarrow$
        multi-core requires global EDF or partitioning, both introducing
        new complexity (migration, cache effects, scheduling anomalies)
  \item \textbf{No resource sharing support:}
        basic CBS does not handle shared resources or priority inversion
        $\Rightarrow$ must be extended (e.g., combined with SRP or PCP)
  \item \textbf{Budget timer overhead:}
        requires one hardware timer interrupt per active server period
        $\Rightarrow$ problematic on systems with limited timer hardware
\end{itemize}

% ============================================================
\section{Hidden or Unrealistic Assumptions}
% ============================================================

\begin{itemize}
  \item \textbf{Assumption: implicit deadlines ($D_i = T_i$):}
        CBS schedulability proof assumes periodic tasks have implicit
        deadlines. Constrained deadlines ($D_i < T_i$) require
        additional analysis.
  \item \textbf{Assumption: fully preemptive system:}
        non-preemptive or limited-preemption RTOS violates CBS
        correctness $\Rightarrow$ blocking time must be accounted for
  \item \textbf{Assumption: reliable WCET estimation:}
        $Q_s$ is a design parameter, not a measured WCET. If aperiodic
        tasks consistently exceed $Q_s$, temporal isolation holds but
        response time degrades significantly
  \item \textbf{Assumption: single aperiodic server:}
        multiple CBS servers on one processor are possible but require
        careful bandwidth allocation to avoid over-commitment:
        $\sum U_{s,k} + \sum U_i \leq 1$
  \item \textbf{Assumption: ideal EDF runtime:}
        EDF requires $O(\log n)$ priority queue management at each
        scheduling event $\Rightarrow$ overhead non-trivial on very
        small MCUs with limited clock speed
\end{itemize}

% ============================================================
\section{Scalability Analysis}
% ============================================================

\begin{itemize}
  \item \textbf{Runtime scalability:}
        server state update $O(1)$ per event $\Rightarrow$ scales well
        with number of aperiodic jobs
  \item \textbf{System size scalability:}
        each additional CBS server adds $U_{s,k}$ to EDF utilization
        $\Rightarrow$ total bandwidth still bounded by Eq.~(1);
        multiple servers are feasible if total $U \leq 1$
  \item \textbf{Memory scalability:}
        constant per server ($Q_s$, $T_s$, $c_s$, $d_s$, queue pointer)
        $\Rightarrow$ scales to many servers on embedded hardware
  \item \textbf{Multi-core scalability:}
        \textbf{does not scale} without major extensions. Global EDF
        with CBS on $m$ cores has open research problems (migration
        overhead, processor affinity, cache pollution)
  \item \textbf{Overload scalability:}
        under sustained overload (aperiodic demand $> U_s$), jobs queue
        and response time grows unboundedly $\Rightarrow$ CBS provides
        isolation but not bounded response under permanent overload
\end{itemize}

% ============================================================
\section{Suitable Application Scenarios}
% ============================================================

\textbf{Scenario 1: Robotic arm controller}
\begin{itemize}
  \item Periodic: joint control loop at 1~kHz (hard RT)
  \item Aperiodic: operator commands, sensor fault events (soft RT)
  \item CBS assigns $U_s$ to command handler; control loop unaffected
        under any aperiodic load
  \item Why suitable: single-core EDF system, soft aperiodic deadlines,
        need for formal bandwidth separation
\end{itemize}

\textbf{Scenario 2: Multimedia streaming on embedded Linux}
\begin{itemize}
  \item Periodic: audio/video decode tasks (hard RT)
  \item Aperiodic: network packet handler (soft RT)
  \item Linux \texttt{SCHED\_DEADLINE} implements CBS-like mechanism
  \item Why suitable: EDF available, temporal isolation needed,
        network bursts must not affect decode quality
\end{itemize}

% ============================================================
\section{Unsuitable Application Scenarios}
% ============================================================

\textbf{Scenario 1: AUTOSAR automotive ECU}
\begin{itemize}
  \item AUTOSAR OS uses fixed-priority scheduling (OSEK standard)
  \item CBS cannot be applied without replacing the scheduler
  \item Why unsuitable: fixed-priority only; use Sporadic Server instead
\end{itemize}

\textbf{Scenario 2: Hard real-time aperiodic tasks}
\begin{itemize}
  \item Example: safety-critical interrupt handler with hard deadline
  \item CBS provides bandwidth guarantee, NOT per-job deadline guarantee
  \item If aperiodic job must meet a hard deadline, CBS cannot certify it
  \item Why unsuitable: no formal per-job response time bound in basic CBS
\end{itemize}

\textbf{Scenario 3: Multi-core SoC with shared cache}
\begin{itemize}
  \item CBS defined for single processor
  \item Multi-core introduces cache contention and migration overhead
        not modeled in CBS
  \item Why unsuitable: formal guarantees break under multi-core
        without specialized extensions
\end{itemize}

% ============================================================
\section{Proposed Extensions and Open Questions}
% ============================================================

\begin{itemize}
  \item \textbf{Extension 1 — Adaptive CBS:}
        dynamically adjust $(Q_s, T_s)$ based on observed aperiodic
        load $\Rightarrow$ better bandwidth efficiency under variable load.
        Open question: how to guarantee stability during parameter transitions?

  \item \textbf{Extension 2 — Multi-core CBS:}
        extend budget/deadline rules to global EDF on $m$ processors.
        Open question: how to handle task migration without violating
        temporal isolation?

  \item \textbf{Extension 3 — CBS with resource sharing:}
        combine CBS with Stack Resource Policy (SRP) or Priority Ceiling
        Protocol (PCP) to handle shared resources.
        Open question: does the schedulability bound change when
        blocking time is added?

  \item \textbf{Extension 4 — Hard CBS variant:}
        Buttazzo describes a stricter replenishment rule (``hard CBS'').
        Open question: does hard CBS improve isolation at the cost of
        response time, and under what workload does it outperform
        standard CBS?

  \item \textbf{Open research question:}
        Is there a closed-form worst-case response time bound for
        aperiodic jobs under CBS (analogous to the DS response time
        analysis under RM)?
\end{itemize}

% ============================================================
\section{Milestone 3 Checklist Mapping}
% ============================================================

\begin{table}[h]
\centering
\caption{M3 Checklist Mapping}
\renewcommand{\arraystretch}{1.3}
\begin{tabular}{|p{5.5cm}|p{2.0cm}|}
\hline
\textbf{M3 Requirement} & \textbf{Section} \\
\hline
Concrete strengths of CBS & Section 1 \\
\hline
Concrete limitations and risks of CBS & Section 2 \\
\hline
Hidden or unrealistic assumptions in CBS & Section 3 \\
\hline
Runtime, memory, scalability implications & Section 4 \\
\hline
Suitable application example for CBS & Section 5 \\
\hline
Unsuitable application example for CBS & Section 6 \\
\hline
Proposed extensions and open questions & Section 7 \\
\hline
\end{tabular}
\end{table}

% ============================================================
\section*{References}
% ============================================================

\begin{thebibliography}{9}

\bibitem{buttazzo}
G.~C.~Buttazzo,
\textit{Hard Real-Time Computing Systems: Predictable Scheduling
Algorithms and Applications}, 4th~ed.
Springer, 2024.
[Primary reference. CBS in chapter on aperiodic task servers.]

\bibitem{abeni1998}
L.~Abeni and G.~Buttazzo,
``Integrating Multimedia Applications in Hard Real-Time Systems,''
in \textit{Proc. 19th IEEE Real-Time Systems Symposium
(Cat. No. 98CB36279)}, IEEE, 1998, pp.~4--13.
[Original CBS paper: formal definition, budget/deadline rules,
schedulability proof.]

\bibitem{abeni2004}
L.~Abeni and G.~Buttazzo,
``Resource Reservation in Dynamic Real-Time Systems,''
\textit{Real-Time Systems}, vol.~27, no.~2, pp.~123--167, 2004.
[Extended CBS analysis: response time results, comparison with TBS.]

\bibitem{ai_usage_m3}
S.~Ferdous,
\textit{AI Usage Protocol, Milestone 3: Constant Bandwidth Server},
HSHL, 2026.
[Protocol v0.2; tool: Claude claude-sonnet-4-6; JSON+bib submitted.]

\end{thebibliography}

\end{document}
